%
%  GALLEGO USAR (LINEA 10)
%  CASTELLANO USAR (LINEA 11)
%  ENGLISH USE (LINE 12)
%
%  MAC (LINEA 13)
%  WINDOWS & LINUX (LINEA 14)

\documentclass[12 pt, a4paper,twoside]{article} %E aceptable tamen 11pt
\usepackage[galician]{babel}\decimalpoint  %Para Galego
%\usepackage[spanish]{babel}\decimalpoint   %Para Castelan
%\usepackage[english]{babel}\decimalpoint    %Para Ingles
\usepackage[latin1]{inputenc} %acentos e caracteres para Mac
%\usepackage[utf8]{inputenc} %acentos e caracteres para Windows & Linux
%ou en vez dos dous renglons de arriba tamen pode usarse ASCII puro e buscar Keyboard Layouts en internet
\usepackage[T1]{fontenc}
\usepackage{bm}
\usepackage{graphicx}
\usepackage{multirow}
\usepackage[table,xcdraw]{xcolor}
\usepackage{hhline}
\usepackage{amsmath, amsfonts, latexsym, color, textcomp, anysize}
\usepackage[rightcaption]{sidecap}
\usepackage{url}
\usepackage{tablefootnote}
\usepackage{fancyhdr}
\usepackage{blindtext}
\usepackage{scrextend}
\addtokomafont{labelinglabel}{\sffamily}
\usepackage{caption}% descomentar de querer texto mais pequeno nas captions de figura e taboas
%\captionsetup{font=footnotesize} % descomentar de querer texto mais pequeno nas captions de figura e taboas. Existen mais opcions en internet.
% \usepackage{subcaption} % eu (sebas) uso sempre este paquete para as figuras e tal, sen meter o caption normal
\usepackage{subfigure}
\parindent=5mm %sangr�a
\parskip=3mm   %separacion entre parrafos
\usepackage{gensymb}
\usepackage{wrapfig}
%\usepackage{natbib}
\usepackage[hidelinks, pdftex]{hyperref}


\definecolor{usc}{RGB}{15,41,118}
\definecolor{grise}{RGB}{234, 236, 240}
\newcommand{\abs}[1]{\lvert#1\rvert}


\begin{document}
\pagenumbering{Roman}
\pagestyle{empty}

\newcommand\TituloDoTraballo{T\'{\i}tulo do Traballo aqu\'{\i}}  % <-- Titulo tal como vai aparecer na portada externa e na interna. Debe coincidir co titulo subido a secretaria virtual no momento do deposito da memoria/solicitude de defensa. Non ten por que coincidir co titulo do proxecto do traballo (o titor debe incluir no seu informe o seu acordo co titulo final). O mais adecuado e que o titulo na portada este no idioma do texto principal da memoria. Tamen pode cambiarse de idioma do resto do contido da portada.

\newcommand\EspecialidadeMaster{Especialidade (valeirar este campo para TFG, cumplimentar para TFM)}  % <-- Deixar valeiro no caso de TFG, e decir introducir {}. No caso de TFM, e IMPORTANTE que conste: Esta considerado un requisito para obter a especialidade da titulacion do master que a devandita especialidade conste no TFM,  sucedendo que na implementacion informatica da USC ao facer a solicitude de defensa de TFM na secretaria virtual a unica opcion onde introducir a especialidade e no propio campo do titulo do traballo. Por eso, a mencion de especialidade debe facerse como unha addenda ao titulo de TFM, a modo de subtitulo - en particular, na secretaria virtual, este subtitulo ou mencion a especialidade debe incluirse no mesmo campo informatico do titulo de TFM (por exemplo introducir no campo titulo do TFM "Estudio de superconductores ultrarelativistas. Especialidad en Fisica Fundamental." Asi vai aparecer tamen no expediente academico do estudante.)

\newcommand\DataDefensa{Febreiro 2020} %<-Mes e ano da presentacion

\begin{center}
%{\sc\large\color{usc} \sc Universidade de Santiago de Compostela}
\vspace{3em}
\includegraphics[width=10em]{USC.png}
\hspace{1cm}
\begin{tabular}[b]{c}
{\large\color{usc} \sc Facultade de F\'{\i}sica} \vspace{0.5em}\\
{\large\color{usc} \sc Grao en F\'{\i}sica } \vspace{0.5em}\\  %<-- Ou Master
{\large\color{usc}  curso 2020-21} \vspace{0.5em}\\%Actualizar curso!
{\Large\color{usc} \sc Traballo de Fin de Grao} %<-- Ou Master
\end{tabular}



\vspace{3cm}
\rule{65mm}{0.2mm}\\
\vspace{1cm}

{\sc\LARGE \TituloDoTraballo}

{\sl\large \EspecialidadeMaster}



\vspace{0.5cm}
\rule{65mm}{0.2mm}\\
\vspace{2cm}
\end{center}


\begin{tabular}{l}
{\sl\large Autor:} \\
{\bf\Large Nome do autor} % nome e apelidos do estudante
\vspace{1em}\mbox{} \\
%
{\sl\large Titor:} \\
{\bf\large Nome do titor} \\
{\sl\large \'Area do titor, Departamento do titor}
\vspace{1em}\mbox{} \\
%
{\sl\large Cotitor: (opcional)} \\
{\bf\large Nome do cotitor} \\
{\sl\large Affiliation do cotitor}
\end{tabular}



\mbox{}
\mbox{}\hfill{\large \DataDefensa} \vspace{1cm}

\noindent{\tiny O autor autoriza a consulta e emprego desta memoria para uso acad\'emico e de investigaci\'on (autorizaci\'on detallada nas p\'axinas interiores).}



\clearpage
\mbox{}
\clearpage %fin portada externa, inicio portada interna%%%%%%%%%%%

\mbox{}\\
Facultade de F\'{\i}sica\\
Grao en F\'{\i}sica \\  
curso 2020-21\\%Nota: actualizar curso
{\sc Traballo de Fin de Grao}\vspace{3cm}\\
%
{\sc\LARGE \TituloDoTraballo}\vspace{1cm}\\
{\sl\large \EspecialidadeMaster}\vspace{2cm}\\
%
{\sl Autor:} {\bf Nome do autor}\\ % nome e apelidos do estudante
{\sl Titor:} {\bf Nome do titor}, {\sl Affiliation do titor, quiz\'ais m\'ais detallada que na portada: grupo investigaci\'on, email, etc}\\
{\sl Cotitor (opcional):} {\bf Nome do titor}, {\sl Affiliation do titor}\vspace{1cm}\\
%
Opcionalmente, inclu\'{\i}r aqu\'{\i}  li\~nas de agradecementos a outras persoas que colaborasen, referencias a financiaci\'on econ\'omica, notas sobre patentes ou publicaci\'ons de traballos derivados do traballo, ou calqueira outra informaci\'on complementaria de contextualizaci\'on xeral do traballo que considere convinte.

\mbox{}

\mbox{}\hfill{Data de presentaci\'on: \DataDefensa}


\clearpage





\thispagestyle{empty}
\pagebreak

{\sl Declaraci\'on asinada polo autor da orixinalidade do traballo}

O autor do traballo declara que o presente \'e un traballo orixinal. Autoriza asemade ao control por persoal da Universidade de Santiago de Compostela da devandita orixinalidade, eventualmente mediante o emprego de bases de datos e a inclusi\'on nelas.

En Santiago de Compostela, a X de Febreiro de 2020. Asinado,\vspace{3cm}


{\sl Autorizaci\'on do autor \'a difusi\'on do traballo (opcional)}

O  autor autoriza \'a difusi\'on do traballo aos efectos considerados nos vixentes regulamentos de TFG e TFM da Universidade de Santiago de Compostela (Artigo 11.3) e de TFM do M\'aster en F\'{\i}sica (Artigo 33), entendendo que esta autorizaci\'on non infl\'ue na propiedade intelectual do traballo nen \'a posibilidade de publicar o mesmo total ou parcialmente por outros medios. Autoriza asemade a que a Facultade de F\'{\i}sica desa Universidade dispo\~na de copia electr\'onica do traballo para o seu arquivo, consulta e emprego para usos acad\'emicos e de investigaci\'on coa menci\'on espec\'{\i}fica ao autor. 

En Santiago de Compostela, a X de Febreiro de 2020. Asinado,\

\thispagestyle{empty}
\pagebreak

%%%%%%%%%%%% En caso de querer incluir unha paxina en branco despois da portada
%\thispagestyle{empty}
%\newpage
%$\ $
%\newpage
\clearpage
\thispagestyle{empty} %%<- poner en mas paginas si los resumenes ocupan varias (elimina la numeracion de la pagina actual)
\pagebreak


\begin{flushleft} \hyphenpenalty=10000\exhyphenpenalty=10000 {\bf $\bullet$ Resumo:\;\;}
%
Aqu\'{\i} vai o resumo en galego. Debe coincidir co introducido na secretar\'{\i}a virtual no momento do dep\'osito da memoria final e solicitude de defensa. Dado que a aplicaci\'on inform\'atica de secretar\'{\i}a virtual non admite calqueira caracter, o regulamento permite introducir nela representaci\'ons alternativas dos caracteres problem\'aticos (por exemplo introducir gamma en vez de $\gamma$, introducir YBa2Cu3O(7-delta) en vez de \mbox{YBa$_2$Cu$_3$O$_{7-\delta}$}, etc.). Non ten por qu\'e coincidir co resumo do proxecto de TFG/TFM que foi feito no momento da asignaci\'on de traballo e titor (entendendo que o titor da o visto bo no seu informe final). Ten que haber resumos en (como m\'{\i}nimo) galego, castel\'an e ingl\'es, cada un de 300 palabras m\'aximo.
%
\end{flushleft}\mbox{}


\begin{flushleft} \hyphenpenalty=10000\exhyphenpenalty=10000 {\bf $\bullet$ Resumen:\;\;}
%
Aqu\'{\i} va el resumen en castellano. El orden de los idiomas puede cambiarse a voluntad. Tambi\'en (y aunque la hoja de res\'umenes no contabiliza para el n\'umero l\'{\i}mite de p\'aginas, ver secci\'on \ref{quecontabiliza} de este documento) puede  reducirse, si se desea, el tama\~no de letra de los res\'umenes en los dos idiomas que no sean el usado  en el texto principal de la memoria (por ejemplo, anteponiendo \verb \footnotesize{}  al texto). % anteponiendo \footnotesize{} al texto, por ejemplo
%
\end{flushleft}\mbox{}


\begin{flushleft} \hyphenpenalty=10000\exhyphenpenalty=10000 {\bf $\bullet$ Abstract:\;\;}
%
The abstract goes here.
%
\end{flushleft}\mbox{}

\begin{center}
\rule{61mm}{0.1mm}\\
\end{center}
{\bf Agradecementos (opcional)}
Aqu� p�dense incluir os agradecementos, se non foron inclu\'{\i}dos xa na portada interior.
 % paxina de resumo e agradecementos. D�bese ter en conta o idioma seleccionado para o documento

\clearpage
\pagestyle{fancy}
\fancypagestyle{plain}
\lhead{}
\chead{}
\rhead{}
\renewcommand{\headrulewidth}{0.1pt}
\lfoot{} 
\cfoot{\thepage}
\rfoot{} 
\renewcommand{\footrulewidth}{0pt}

\pagenumbering{arabic}\setcounter{page}{1}
\tableofcontents
\clearpage

%%%%%%% Incluimos as secci�ns que sexan necesarias
\section{Introducci\'on}

Aqu\'{\i} ven o texto da introducci\'on. 

Este documento \'e unha plantilla para elaborar memorias de TFG e TFM en F\'{\i}sica. {\bf O USO DE ESTA PLANTILLA NON \'E OBRIGATORIO}. \'E un documento en \LaTeX. O seu c\'odigo (nos ficheiros~ .tex) \'e bastante  auto-explicativo do seu uso. A\'{\i}nda as\'{\i}, o texto que empregamos como exemplo tam\'en resposta a algunhas preguntas frecuentes.

Nesta plantilla o texto vai escrib\'{\i}ndose en varios ficheiros, pero o documento principal a compilar \'e {\tt tfg.pdf}, que chama a todos os dem\'ais. Lembra  grabar todos os ficheiros nos que traballes!

Esta plantilla pode modificarse e/ou distribuirse a vontade, sen necesidade de permiso. 
\clearpage
\section{Algunhas Preguntas Frecuentes}

Neste cap\'{\i}tulo inclu\'{\i}mos algunha subsecci\'on, ecuaci\'on, figura e t\'aboa, e de paso respostamos a algunha pregunta frecuente.


\subsection{?`Qu\'e contabiliza para o n\'umero m\'aximo de p\'axinas l\'{\i}mite (25  m\'aximo para TFG e 40 para TFM)?}\label{quecontabiliza}
 S\'{\i} contabilizan o \'{\i}ndice, a introducci\'on e resto de cap\'{\i}tulos, a bibliograf\'{\i}a, conclusi\'ons, figuras e t\'aboas, as\'{\i} como calqueira ap\'endice ou listado de c\'odigo inform\'atico. \'E dicir, b\'asicamente dende o \'{\i}ndice (inclu\'{\i}do) ata o final excepto as contraportadas internas e externas.
Non contabilizan a portada externa nin interna, as declaraci\'ons de orixinalidade e autorizaci\'on de difusi\'on asinadas, nin os  res\'umenes, nin as contraportadas internas e/ou externas meramente protectoras ou con informaci\'on de peche.

\subsection{?`C\'omo introduzo acentos, etc.?}
Vai \'as primeiras li\~nas de c\'odigo do ficheiro principal ({\tt tfg.tex}) e ver\'as c\'omo empregar os acentos ``normais'' para varios sistemas operativos. Alternativamente, nos nosos textos de exemplo ver\'as c\'omo escribir en ASCII simple, independente do sistema operativo, tanto acentos coma outros caracteres como por exemplo \~n, \"u, !`, ?`, \%, @, hbar. En internet hai moita informaci\'on (busca latex keyboard layouts, latex symbols, etc.).

\subsection{?`C\'omo fago unha ecuaci\'on?}
No medio do texto ser\'{\i}a $\lambda=2\pi/\hbar$. Nunha li\~na propia ser\'{\i}a  (mira no c\'odigo aqu\'{\i}): 
% 
\begin{equation}
\lambda=\frac{2\pi}{\hbar}.
\label{eq-lambda}
%\nonumber se non queremos numero de ecuacion (non se recomenda)
\end{equation}
% se deixas un renglon en blanco no codigo comezas un parrafo novo. Para evitalo podes usar un renglon cun caracter de comentario %
Ou incluso
%
\begin{subequations}\label{todaaecuacion}%este entorno permite numerar ecuacions como 2a, 2b etc. O label e opcional para se queres chamar logo a ecuacion
\begin{align}%este entorno sirve para que os signos igual esten ordeados
\lambda_1&=\frac{2\pi}{\hbar},\\
\lambda_2&=\frac{4\pi}{\hbar}+\int_{-\infty}^0{\rm d}x \sqrt{\sum_i\omega(x)}.\label{partedaecuacion}
\end{align}
\end{subequations}
%

\subsection{?`C\'omo referencio a unha secci\'on, ecuaci\'on, referencia bibliogr\'afica, etc.?}
Tes que deixar nomeadas as {\tt label} ou etiquetas no obxecto que queres referir, e logo podes chamalas (mira o c\'odigo desta li\~na para ver a sintaxe; mira tam\'en o c\'odigo no obxecto referido) as\'{\i}: ``A secci\'on~\ref{quecontabiliza}''. ``As ecuaci\'ons{\ref{eq-lambda}, \ref{todaaecuacion} ou~\ref{partedaecuacion}''. ``O cadro~\ref{taboa1}''. ``a figura~\ref{figura-exemplo}''. En varias implentaci\'ons de \LaTeX\ tes que compilar d\'uas veces o documento cada vez que cambias unha referencia.

%As referencias bibliogr\'aficas def\'{\i}nense en {\tt bibliografia.tex} como ``bibitems'', que definen unha etiqueta ao que chamar no texto, como facemos nesta frase.\cite{patricia} Tam\'en poden ser notas textuais.\cite{nota}

\subsection{T\'aboas, cadros, etc.}
As t\'aboas computan para o l\'{\i}mite m\'aximo de p\'axinas da memoria. Un exemplo \'e o cadro~\ref{taboa1}.
%
\begin{table}[h]
\begin{center}\begin{tabular}{|c|c|r|}
\hline
mostra & masa ($\mu$g) & cor\\
\hline 
mostra 1 & 0.2 & vermella\\
mostra 2 & 0.1 & azul\\
\hline
\end{tabular}\end{center}
\caption{\sl Exemplo de cadro. Neste campo de  {\tt caption} tam\'en inclu\'{\i}mos un tipo de letra diferente, o cal \'e opcional. No pre\'ambulo do ficheiro  {\tt tfg.pdf} hai comentada algunha li\~na con opci\'ons m\'ais potentes, como cambiar o tama\~no de letra nas captions.}\label{taboa1}
\end{table}
%

\subsection{Figuras}
As figuras tam\'en computan para o l\'{\i}mite m\'aximo de p\'axinas. Existen varias formas de inclu\'{\i}r ficheiros de figuras, que podes buscar en internet. Podes  xerar a figura en calqueira programa, imprimila cunha impresora virtual a~.pdf ou outro formato, e chamar no c\'odigo \LaTeX\ ao ficheiro como aqu\'{\i} facemos por exemplo co logo da USC. Respecto a tama\~nos, cuida que o caracter m\'ais pequeno (inclu\'{\i}ndo sub\'{\i}ndices etc.) non te\~na menos altura base (tama\~no dun x) que 2~mm unha vez reducido ao seu tama\~no  no documento final.
%
\begin{figure}[htbp]
\begin{center}
  \includegraphics[width=0.2\linewidth]{USC.png}
  \mbox{}\hskip1cm\mbox{}
  \includegraphics[width=0.2\linewidth]{USC.png}
\caption{Aqu\'{\i} incrustamos d\'uas figuras pero obviamente pode ser s\'o unha, ou m\'ais de d\'uas, regulando os {\tt width} do c\'odigo. Se son varias, recom\'endase etiquetalas como a), b), etc., no propio ficheiro gr\'afico (o noso logo de ejemplo no o fai, claro) e explicar cada unha neste campo {\tt caption}: A subfigura a) amosa blabla.... A subfigura b) amosa bleble. }\label{figura-exemplo}
\end{center}
\end{figure}}

\clearpage
\section{Nome da secci�n}
Contido da secci�n
\cleardoublepage %cleardoublepage pode meter paxinas en branco. Non e obrigatorio. Tampouco para a biliografia, ver abaixo. Se en cambio queres que non haxa cambios de paxina antes dunha seccion, mete os ficheiros con \input{filename} en vez de con \include (poderia precisar a ruta e a extension do ficheiro)
\section{Nome da secci�n}
Contido da secci�n
\cleardoublepage
\section{Conclusi�ns}
Apartado para as conclusi�ns.

\'E habitual que as memorias de TFG e TFM rematen cunha secci\'on de conclusi\'ons. Pero non existe a obriga de que esa secci\'on sexa le\'{\i}da literalmente na exposici\'on oral (obriga-costume que existe nas lecturas de tesis doutorais). Nas defensas orais de TFG e TFM, \'e habitual rematar cun resumo das conclusi\'ons, pero non ten por qu\'e ser unha lectura literal desta secci\'on. Esta secci\'on contabiliza na contaxe de p\'axinas de cara ao l\'{\i}mite m\'aximo da extensi\'on da memoria, as\'{\i} como a secci\'on de bilbliograf\'{\i}a.


%%%%%%%Bibliograf�a
\appendix
\cleardoublepage
\markboth{REFERENCIAS}{REFERENCIAS}  
\addcontentsline{toc}{section}{Referencias}
\renewcommand{\refname}{Referencias}  %<--cambiar ese renglon e os de arriba segun o idioma empregado. Pode chamarse tamen bibiografia no caso de que non haxa notas coma o do bibitem{nota} do noso exemplo

\begin{thebibliography}{99}
%\footnotesize %descomentar para cambiar o tamano de letra desta seccion (footnotesize e valido)
%Exemplo
\bibitem{patricia} J. Abraham et al., ``The fluorescence detector of the Pierre Auger Observatory'', {\it Nuclear Instruments and Methods in Physics Research A}, vol. 620, pp. 227-251, 2010.
\bibitem{nota} \'As veces as entradas bibliogr\'aficas poden conter textos explicativos, como se fosen notas ao p\'e. Tam\'en poden ir varias referencias bibliogr\'aficas nunha soa entrada. Non \'e obrigatorio ning\'un formato espec\'{\i}fico de referencias (orde dos volume-p\'axinas-ano, por exemplo) pos hai moitos en uso internacionalmente. A secci\'on de bibliograf\'{\i}a computa na contaxe de p\'axinas de cara ao l\'{\i}mite m\'aximo da extensi\'on da memoria. No c\'odigo do ficheiro {\tt bibliografia.tex} pode usarse a opci\'on de cambiar tama\~nos, e incluir outras. Lembra compilar varias veces o documento final para que as referencias est\'en correctamente actualizadas.
\end{thebibliography}


\end{document}
